% !TEX root = ../main.tex
\chapter{서론과 연구 문제}
\label{ch:introduction}

디파이(DeFi, 탈중앙 금융)는 누구나 볼 수 있는 블록체인 위에서 돈이 오가는 방식을 크게 바꾸었어요. 이제는 은행 같은 중간 기관 없이도 돈을 빌려주고, 빌리고, 빌린 돈을 보태 거래할 수 있어요\cend{1}. 그런데 신용 거래 비슷한 일은 아직도 대부분 담보를 많이 맡겨야만 할 수 있어요. 디파이 프로그램은 블록체인 밖의 신분 확인 제도나 이미 쌓인 신용 기록을 쓸 수 없기 때문이에요\cend{2}. 진짜 이름 대신 주소로만 활동하는 세상(가명성)에서는, 서로 모르는 사람들 사이에 돈을 어떻게 잘 맡길지가 가장 큰 문제예요.

이 논문은 ERC-20(토큰 공통 규칙) 허락(allowance, 토큰을 얼마까지 대신 써도 된다는 허락) 화살표(엣지)가 페이지랭크(PageRank)에 무엇을 보태는지 물어요. 여기서 페이지랭크는 온체인(블록체인 위의) 지갑의 평판을 재는 잣대예요. 허락 화살표는 \texttt{approve}와 \texttt{permit}이 남기는 \texttt{Approval} 이벤트(계약이 남기는 기록)에서 나와요. 엔도스랭크(EndorseRank)는 이 화살표 위에서 계산한 페이지랭크예요. 화살표 무게(가중치)는 Do, Do, Nguyen\cend{3}의 AWP(적응형 가중 페이지랭크)와 똑같이 매겨요. AWP는 송금(transfer)을 바탕으로 한 방법이에요. 두 점수는 같은 아비트럼 원(Arbitrum One) 지갑들에서 계산해요. 이때 두 종류의 화살표를 함께 걷는 결합 걷기와, 걷기들이 매끄럽게 다듬는(평활하는) 원시 차수(화살표를 그냥 센 수)도 함께 계산해요. 검증은 AWP에서 쓴 순위 상관(두 줄 세우기가 얼마나 비슷한지 재기) 방법을 따라요\cend{3}. 여기에 짝지은 대비(같은 지갑들에서 두 점수를 맞대어 본 차이), 원시 차수 기준선(비교 기준), 시간 홀드아웃(나중 기간을 떼어 두고 맞혀 보는 시험)을 더했어요. 이 논문에서 온체인 평판 지표는 온체인 활동과 주고받은 일에서 끌어낸 값이에요. 오프체인(블록체인 밖의) 신분 정보는 하나도 쓰지 않아요\cend{4}.

\dissertationSubheading[sec:background]{배경: 디파이, 담보 맡기기, 가명성}

누구나 볼 수 있는 블록체인에서는 중앙 기관 없이도 프로그램으로 거래를 마무리(결제)할 수 있어요\cend{5}. 스마트 계약(블록체인 위에서 저절로 실행되는 프로그램)은 거래 조건을 코드로 담을 수 있어요\cend{6}. 이런 계약을 쓰는 디파이 플랫폼은 거래를 돈을 맡아 주는 곳(수탁 기관)에 보내지 않고 블록체인 위에서 바로 처리해요\cend{7}. 디파이 서비스들은 이런 계약을 겹겹이 쌓아 빌려주기(대출), 자동 시장 조성(프로그램이 늘 사고팔 값을 내놓는 일), 무기한 선물(정해진 끝 날짜 없이 앞으로의 값에 거는 거래)을 위한 탈중앙 시장을 만들어요\cend{8}. 아비트럼 원과 다른 레이어 2(L2) 네트워크(바탕 블록체인 위에서 돌아가는 블록체인)에서는 수수료가 아주 싸서 보통 사용자도 자주 거래해요. 그리고 모든 송금, 허락, 서비스 이벤트를 지갑마다 하나하나 살펴볼 수 있어요.

이렇게 다 보이는데도, 디파이의 신용 시장은 본래 과담보(빌리는 돈보다 담보를 더 맡기는) 방식에 머물러 있어요. 대부분의 디파이 서비스는 빌리는 사람이나 트레이더(거래하는 사람)에게 포지션(지금 열어 둔 거래)의 값어치보다 더 많은 담보를 맡기게 해요. 거래 상대가 가명의 계정이고, 그 주인이 누구인지 확인되지 않았기 때문이에요\cend{9,10}. 가명성은 누구나 볼 수 있는 장부의 기본 성질이에요. 누구든 거의 돈을 들이지 않고 주소를 만들 수 있고, 한 주체(사람이나 단체)가 지갑을 여러 개 가질 수 있어요\cend{11}. 보통의 신용 등급은 개인이 지나온 기록에 기대는데, 그런 기록은 이 환경에 잘 옮겨지지 않아요\cend{2}. 그래서 디파이 서비스는 개인별 평판 대신 담보로 쌓아 둔 돈, 청산(손실이 너무 커져 거래가 강제로 끝나는 일) 장치, 서비스 규칙으로 위험을 다스려요.

담보는 디파이 서비스를 지켜 주지만, 돈을 알뜰하게 쓰지 못하게 해요. 트레이더와 빌리는 사람은 다른 곳에 투자할 수도 있었던 돈을 묶어 두어야 해요. 또 전통 금융은 신용 기록이 있는 사람에게 담보를 덜 맡겨도 되는 상품을 내주지만, 디파이 플랫폼에는 그런 상품이 하나도 없어요. 그래서 연구자들은 오프체인 신용 기록을 대신할 온체인 신호로 눈을 돌렸어요. 송금 그래프(점과 화살표로 이은 그림)의 중심성(그래프 안에서 얼마나 중심에 있는지)\cend{4}, 네트워크를 따라 위험 퍼뜨리기(전파)\cend{12}, 전망 이론(사람이 이익과 손해를 어떻게 느끼는지 다루는 이론)과 블록체인을 바탕으로 한 신용 점수 매기기\cend{13}, 대출 서비스 자료를 다루는 머신러닝(기계 학습) 처리 과정(파이프라인)\cend{14} 같은 것들이에요. 이 연구들은 저마다 같은 가정을 해요. 공개 장부 자료에는 사용자가 누구인지 밝혀내지 않고도 위험을 판단하는 데 쓸 수 있는 정보가 들어 있다는 가정이에요.

그래프 순위 알고리즘(정해진 순서대로 하는 계산 방법)은 이 문제에 잘 맞아 보여요. 페이지랭크와 그 변형들은 되풀이해 이어지는 보증(믿고 맡김)으로 점(노드)에 순위를 매겨요. 그래서 한 점의 점수는 다른 점들에게서 받은 보증에 따라 정해져요\cend{15}. EigenTrust\cend{16}와 TrustRank\cend{17} 같은, 그보다 앞선 웹 신뢰 알고리즘들은 비슷한 생각을 P2P(사용자끼리 직접 이어진) 신뢰와 웹 스팸에 썼어요. Do와 Do\cend{4}는 송금 그래프에 페이지랭크와 HodgeRank를 함께 써서 이더리움(Ethereum) 사용자를 얼마나 믿을 수 있는지 판단했어요. 그리고 그 결과를 일종의 사회적 신용이라고 불렀어요. Do, Do, Nguyen\cend{3}이 제안한 적응형 가중 페이지랭크(AWP)도 송금 그래프로 순위를 매기는 알고리즘이에요. AWP는 송금마다 무게를 달아요. 무게는 금액에 대한 상한이 있는 함수(금액이 아무리 커도 일정한 값을 넘지 않게 하는 식)와, 송금이 얼마나 오래되었는지에 대한 로지스틱 함수(오래된 송금일수록 가볍게 세는 식)로 정해요. 또 최근에 송금을 보낸 주소에서 걷기를 다시 출발(재시작)해요. AWP를 만든 사람들은 AWP가 들어온 송금의 수, 금액과 함께 움직인다(상관관계가 있다)는 것을 보여 AWP를 검증했어요. 들어온 송금의 수와 금액은 경제적 중요성의 대리 지표(대신 재는 값)로 쓴 것이에요. 이 연구들 덕분에 송금 그래프 평판은 믿을 만한 기준선이 되었어요. 하지만 이 연구들은 돈을 쓰도록 분명하게 허락하는 일을 빼놓았어요. 이것은 ERC-20 허락으로 기록되는, 따로 떨어진 온체인 행동이에요\cend{18,19}.

\dissertationSubheading[sec:motivation]{온체인 평판이 필요한 까닭}

빌린 돈을 보태 거래하는(레버리지) 디파이에서 온체인 평판 연구가 필요한 까닭은 크게 네 가지예요.
\begin{itemize}
\item 돈을 알뜰하게 쓰기: 디파이 서비스들은 순서를 매기는 신호(누가 앞이고 누가 뒤인지 알려 주는 신호)를 찾고 있어요. 이런 신호는 등급별 서비스를 뒷받침할 수 있어요. 예를 들면 한도를 더 자주 올려 주기, 수수료 낮춰 주기, 감시(모니터링)에서 더 높은 우선순위 주기 같은 것이에요. 어쩌면 앞으로 통제된 실험(조건을 정해 두고 하는 실험)에서는 담보 요구를 지갑마다 다르게 하는 데까지 쓸 수도 있어요. 그러면서도 거래 마무리는 그대로 블록체인 위에서 해요\cend{1}.
\item 따져 볼 수 있음: 오프체인 신용 점수는 속을 들여다볼 수 없어요. 하지만 공개된 이벤트 자료로 계산한 그래프 점수는 누구나 다시 계산하고 따져 볼(감사할) 수 있어요. 그래서 속을 알 수 없는(블랙박스) 탈중앙 신용 점수 매기기에 대한 걱정을 조금은 덜어 줄 수 있어요\cend{2}.
\item 운영 비용: 자주 새로 고치는 점수는 다시 계산하는 비용이 적게 들어야 해요. 예를 들어 레버리지를 바꿀 때마다 그 전에 새로 고치는 점수가 그래요.
\item 표준 만들기: 다른 시장에서는 거래 상대를 가려낼 때 심사용 비율을 써요. 하지만 가명으로 움직이는 디파이에는 지갑 하나하나를 그렇게 가려낼, 표준이 되는 순서 언어가 없어요. \ref{ch:discussion}장에서는 이 논문에서 살펴본 점수들이 어디서 그 역할을 할 수 있고, 어디서 할 수 없는지 이야기해요.
\end{itemize}

송금을 바탕으로 한 AWP는 공개되어 있고 누구나 따져 볼 수 있어요. 하지만 돈을 낼 때마다 그래프가 커져요\cend{3}. 허락 그래프는 짜임새가 다른 대안이에요. 사용자가 ERC-20 토큰을 허락할(approve) 때마다, 계약은 스펜더(spender, 허락을 받아 대신 쓰는 쪽)가 그 토큰을 얼마까지 옮길 수 있는지 기록해요\cend{18,19}. 새 허락은 예전 허락을 덮어써요. 그래서 주인(토큰 주인)--스펜더 쌍과 토큰마다, 마지막 허락(가장 최근 허락)이 그 쌍의 지난 기록 전체를 대신해요. 또 허락은 송금보다 드물어요. 그래서 같은 지갑들에서 보증 그래프 $G_{\text{approve}}$는 송금 그래프 $G_{\text{transfer}}$보다 화살표가 열 배 정도 적어요. 페이지랭크를 푸는 비용은 $|E|$(화살표 수)에 비례해 커지므로, 계산도 그만큼 빨리 끝나요(표~\ref{tab:benchmark-runtime}). 두 경우 모두 같은 풀이 프로그램(solver)을 돌려요. 그래서 속도 향상(몇 배 빠른지)은 더 듬성듬성한 그래프 덕분이지, 알고리즘이 나아져서가 아니에요. 주장~C1은 이것을 그래프를 나타내는 방식이 가진 운영상의 성질로 다뤄요.

효율 말고도, 허락은 뜻도 달라요. 송금은 돈이 거쳐 가는 길(라우팅)이나 시장 조성, 또는 그저 받기만 한 것을 나타낼 수 있어요. 반면 허락은 돈을 써도 된다고 적극적으로 내준 허가예요. 허락을 보증 화살표로 다루면, 페이지랭크가 링크를 인용(추천)처럼 읽는 방식\cend{15}을 디파이 서비스들이 이미 기대고 있는 신뢰 관계로 옮겨 오게 돼요. 이렇게 해서 순위가 달라지는지, 그리고 그 순위가 허락 점검, 송금 점검과 맞는 정도가 서로 다른지를 살펴봐요. 대상은 아비트럼 원의 GMX~V2 거래 지갑 $n=5{,}521$개로 된 매칭 표본이에요.\footnote{매칭이란 모든 방법이 똑같은 지갑들에 점수를 매긴다는 뜻이에요. 그 지갑들은 2025년 12월 1일부터 2026년 5월 31일까지 포지션을 세 번 이상 닫은 GMX~V2 지갑이에요. 이 목록은 허락이나 송금을 하나라도 꺼내기 전에 GMX 자료로 정해 두었어요(\ref{sec:data-sources}절).}

\dissertationSubheading[sec:research-problem]{연구 문제}

적응형 가중 페이지랭크(AWP) 방법은 무게를 단 합과 로지스틱 시간 감쇠(시간이 지날수록 무게를 줄이기)를 써서 송금 기록에서 신뢰를 짐작해요\cend{3}. AWP는 경제적으로 풀이할 거리가 많아요. 하지만 AWP는 (i) 송금 기록 전체를 모아야 하고, (ii) 값어치가 큰 지갑을 들어오는 송금을 많이 받는 지갑으로 보고, (iii) 돈을 쓰도록 허락하는 일에는 아무 역할도 주지 않아요. 우리가 제안하는 엔도스랭크는 AWP의 화살표 무게와 풀이 프로그램을 그대로 써요. 하지만 방향이 있는 그래프는 주인--스펜더 쌍과 토큰마다 마지막 허락으로 만들어요. 그렇게 만든 그래프는 훨씬 듬성듬성하고, 담은 신호도 달라요. 또 관계마다 다른 쪽에 순위를 매겨요. 허락은 주인에서 스펜더 쪽을 가리켜요. 그래서 엔도스랭크는 라우터(거래를 이어 주는 계약)나 금고(vault)처럼 허락을 받는 스펜더에게 점수를 매겨요. 그리고 허락을 하나도 받지 않은 주인에게는 모두 똑같은 값을 줘요.

중심 연구 질문은 이거예요. 송금 걷기, 그리고 두 걷기가 매끄럽게 다듬는 원시 차수와 비교할 때, ERC-20 허락 화살표는 아비트럼 원 지갑의 평판 점수인 페이지랭크에 무엇을 보탤까요? 연구 질문(RQ) RQ1--RQ5(\ref{sec:research-questions}절)가 이 질문을 작게 나눠요. 두 종류의 화살표는 같은 주소들 위에 있어요. 그래서 걷기 하나가 점마다 정해진 섞는 비율 $\lambda$로 두 화살표를 섞을 수도 있어요. 이렇게 정한 것이 결합 페이지랭크(C-PR, 두 그래프를 함께 걷는 점수)예요. $\lambda=1$과 $\lambda=0$에서는 한 층(화살표 한 종류로 만든 그래프)만 걸어요. 이렇게 섞는 것을 시험할 만한 까닭은 분석 단위 때문이에요. 트레이더 대부분은 주인이라서 허락 점수로는 서로 가려낼 수 없어요. 두 층을 한 번에 걷는 걷기라면, 송금 층이 트레이더들을 두루 덮는 장점은 지키면서 스펜더에 대한 허락 신호도 함께 실을 수 있을지 몰라요.

점수들은 표~\ref{tab:primary-methods-preview}에 있어요. 표~\ref{tab:hrq-claim-map}에는 가설(H1--H3), 연구 질문(RQ1--RQ5), 그리고 \ref{ch:results}장에서 나온 실증 주장(C1--C5, 자료로 확인한 주장)을 정리했어요.

\begin{table}[htbp]
\centering
\caption{이 논문에서 평가한 점수들(매칭 코호트: 모든 점수가 똑같이 점수를 매기는 지갑 묶음, $n=5{,}521$).}
\label{tab:primary-methods-preview}
\fitwidth{%
\begin{tabular}{p{0.18\linewidth}p{0.28\linewidth}p{0.48\linewidth}}
\toprule
점수 & 입력 그래프 & 역할 \\
\midrule
EndorseRank & AWP와 같은 방식으로 무게를 단 마지막 ERC-20 허락($G_{\text{approve}}$) &허락으로 권한 맡기기(위임)예요. 이 연구에서 살펴보는 대상이에요 \\
AWP & 송금($G_{\text{transfer}}$) & 송금 흐름이에요. Do, Do, Nguyen \cend{3}이 발표한 송금 그래프 방법이에요(\ref{sub:awp-graph}절) \\
C-PR & 두 층 모두, 금액 그대로, 점마다 비율 $\lambda$로 섞은 것 & 두 종류의 화살표가 합쳐지는지 시험하는 결합 연산자(두 층을 섞어 걷는 계산 규칙)예요(RQ5). $\lambda=1$과 $\lambda=0$에서는 한 층만 걸어요. C-PR의 빼 보기 실험은 S-PR로, 허락 걷기에서 다시 출발하는 송금 걷기예요 \\
원시 차수 & $G_{\text{approve}}$와 $G_{\text{transfer}}$에서 점마다 센 것 & 받은 허락 수(허락 입차수)와 송금 입차수(송금을 보내온 서로 다른 주소 수)예요. 페이지랭크들이 매끄럽게 다듬는, 그 점에서 바로 센 수예요. 같은 기간 안에서는 대리 지표이고, 기간 밖에서는 기준선이에요 \\
\bottomrule
\end{tabular}
}
\end{table}

잘못 읽는 일이 없도록, 여기서 주요 용어의 뜻을 정해 둘게요.

\begin{itemize}
\item \emph{디파이(DeFi, 탈중앙 금융):} 빌려주기와 빌리기 같은 금융 서비스예요. 누구나 볼 수 있는 블록체인 위에 스마트 계약으로 만들어져 있어요. 그리고 대개 돈을 맡아 두거나(보관) 거래를 마무리하는(결제) 일을 맡은 중앙 운영자 없이 돌아가요\cend{1}.
\item \emph{스마트 계약:} 블록체인 위에서 저절로 실행되는 프로그램 코드예요. 믿고 맡길 중간 기관 없이 약속한 내용을 실행해요\cend{6}. 이 논문에서 허락 로그와 송금 로그(계약이 남기는 기록)는 스마트 계약이 남겨요.
\item \emph{지갑 / 주소:} 돈을 담고 스마트 계약과 주고받는 블록체인 계정이에요. 우리가 평판 순위를 매기려는 대상이에요\cend{7}. 지갑은 저마다 주소로 구별해요. 오프체인 신분은 미리 가정하지도 않고, 알아내지도 않아요.
\item \emph{온체인 평판 점수:} 온체인 자료로 계산한 수예요. 가명성 아래에서 금융 활동에 대해 한 지갑이 상대적으로 어디쯤 있는지를 나타내요\cend{4}. 점수는 순서만 나타내고, 순위를 매기는 데 써요. 점수 값의 절대적인 크기는 신용 한도에 대해 아무 뜻이 없어요.
\item \emph{페이지랭크:} 방향이 있는 그래프에서 감쇠 계수(화살표를 따라갈 확률)를 써서 영향력을 퍼뜨리는 그래프 순위 알고리즘이에요\cend{15}. 적응형 가중 페이지랭크(AWP)는 송금 그래프 위의 페이지랭크예요. 화살표 무게는 송금마다 금액과 지난 시간으로 정하고, 최근에 활동한 보낸 쪽(송신자)에서 다시 출발해요\cend{3}. 엔도스랭크는 AWP의 화살표 무게를 마지막 허락 그래프에 쓰고, 모든 점에서 고르게 다시 출발해요(\ref{sec:graph-construction}절).
\item \emph{결합 페이지랭크(C-PR):} 허락 층과 송금 층을 함께 걷는 무작위 걷기 하나예요. 지갑마다, 그 지갑이 들어 있는 층들에 대해 두 줄(이동 행렬의 행)을 무게 $\lambda$와 $1-\lambda$로 섞어요(\ref{sub:coupled-graph}절). 보증 씨앗 페이지랭크(S-PR)는 송금 걷기를 그대로 두지만, 허락 걷기의 결과에서 다시 출발해요. TrustRank\cend{17}가 씨앗 집합(믿을 만하다고 미리 고른 점들)에서 다시 출발하는 것과 같아요.
\item \emph{ERC-20 허락 / 허락하기(Approve) / 서명 허락(Permit):} 스펜더가 주인 대신 주인의 토큰을 얼마나 쓸 수 있는지 기록하는 계약 기능이에요\cend{18,19}. \texttt{Approve}는 보통의 온체인 호출이에요. EIP-2612의 \texttt{permit}은 서명으로 똑같은 상태 변화(계약에 적힌 값이 바뀌는 것)를 만들어요. 엔도스랭크는 주인--스펜더 쌍과 토큰마다 0이 아닌 마지막 허락을 써요.
\item \emph{보증 그래프:} 화살표가 돈을 쓸 권한을 분명하게 맡긴 것(위임)을 나타내는, 방향이 있는 그래프예요. 실제로 자산을 옮긴 것과는 달라요\cend{18}. 화살표 $u \to v$는 지갑 $u$가 지갑 또는 계약 $v$에게 $u$의 토큰을 써도 된다고 허락했다는 뜻이에요.
\item \emph{영역 일치도(domain alignment):} 평판 순위가 관찰할 수 있는 대리 지표의 순위와 얼마나 맞는지를 말해요. 스피어먼의 로($\rho$, 두 순위가 얼마나 비슷한지 나타내는 수)\cend{20}와 켄달의 타우($\tau$, 두 줄 세우기가 얼마나 같은 순서인지 나타내는 수)\cend{21}로 재요. 그리고 Cronbach와 Meehl\cend{22}이 말한 뜻의 구성 타당도(점수가 재려는 것을 정말 재는지)로 풀이해요. 어떤 방법은 자기 구성개념(점수가 재려고 하는 생각)의 그래프와는 잘 맞지만, 다른 점검과는 약하게 맞을 수 있어요.
\end{itemize}

\dissertationSubheading[sec:supplementary-extension]{시간 홀드아웃}

같은 기간의 허락에 대한 켄달의 $\tau$(표~\ref{tab:alignment-hybrid})는 의도한 구성개념 점검(점수가 본디 재려던 것과 맞는지 보는 점검)이에요. $G_{\text{approve}}$ 위의 페이지랭크는 스펜더 쪽의 허락 차수와 맞아야 해요. 이 점수는 그 차수를 그래프 전체에 걸쳐 매끄럽게 다듬은 것이기 때문이에요. 점수가 같은 기간의 그래프 너머의 정보를 전하는지 시험하려고, 점수들을 2026년 2월 28일에 고정했어요(동결). 그리고 2026년 3월부터 5월까지에야 알 수 있는 두 라벨(나중에 채점에 쓰는 정답 값)과 견주었어요. 하나는 새 허락 쌍(새로 생긴 주인--스펜더 허락)으로, 보증 구성개념이에요. 다른 하나는 새 송금자(새로 송금을 보내온 주소)로, 흐름 구성개념이에요. 홀드아웃 코호트(묶음, $n=1{,}335$)는 꺼낸 허락 자료에서 $t_1$(2026년 2월 28일 23:59:59~UTC)에 0보다 큰 마지막 허락을 하나 이상 가진 스펜더 전부예요. 또 각 페이지랭크를 $t_1$에서 그것이 다듬는 원시 차수와도 견주었어요. 그래서 퍼뜨린 점수가 그 자리에서 센 수보다 얼마나 나은지 바로 볼 수 있어요.

두 번째 기간은 이 봄 홀드아웃의 결과를 안 뒤인 2026년 9월 28일에 설계했어요. 그리고 2026년 10월 1일에 등록(계획을 미리 공개해 고정해 두기)했어요. 이 기간의 점수는 2026년 5월 31일에 고정하고, 라벨은 2026년 6월부터 8월까지 세어요. 등록을 커밋(저장소에 기록)했을 때, 이 달들의 로그는 아직 꺼내지 않은 상태였어요. 두 번째 기간은 스펜더 라벨 두 개를 되풀이하고, 매칭된 트레이더를 위한 라벨 하나를 더해요. 바로 트레이더가 GMX에서 포지션을 닫은 것 가운데 청산으로 끝나지 않은 비율이에요(\ref{sub:fresh-holdout}절).

\begin{table}[htbp]
\centering
\caption{가설, 연구 질문, 실증 주장 요약.}
\label{tab:hrq-claim-map}
\fitwidth{%
\begin{tabular}{p{0.12\linewidth}p{0.42\linewidth}p{0.40\linewidth}}
\toprule
항목 & 다루는 내용 & 보고한 곳 \\
\midrule
H1 / RQ1 & EndorseRank와 AWP는 서로 다른 지갑 순서를 만드나요? & \ref{sec:rank-divergence}절; 방법 사이의 $\tau$와 95\% 신뢰구간(믿을 만한 범위); 표~\ref{tab:tie-sensitivity}; 주장 C3 \\
H2 / RQ2 & 각 점수는 허락 점검, 그리고 송금 점검과 얼마나 잘 맞나요? & 표~\ref{tab:alignment-hybrid}와~\ref{tab:tie-sensitivity}; 주장 C2 \\
H3 / RQ3 & 같은 $n$에서 각 점수를 계산하는 데 얼마나 드나요? & 표~\ref{tab:benchmark-runtime}--\ref{tab:benchmark-scaling}; 주장 C1 \\
RQ4 & $t_1$의 점수가 $t_2$의 새 허락과 새 송금자를 예측하나요(미리 맞히나요)? 그리고 어떤 페이지랭크라도 자기가 다듬는 원시 차수보다 나은가요? & \ref{sec:holdout-results}절; 표~\ref{tab:holdout-spenders}, \ref{tab:holdout-diff}, \ref{tab:holdout-receiving}, \ref{tab:holdout-traders}와~\ref{tab:fresh-posthoc}; 주장 C4 \\
RQ5 & 결합 연산자 C-PR은 송금만 걷는 걷기보다 나은가요? 그리고 어느 한 층만 걷는 걷기보다 지갑 순위를 더 잘 매기나요? & \ref{sec:coupled-results}절, \ref{sec:holdout-results}절, \ref{sub:fresh-results}절; 표~\ref{tab:alignment-hybrid}, \ref{tab:tau-diff-hybrid}, \ref{tab:holdout-diff}, \ref{tab:fresh-holdout}와~\ref{tab:fresh-contrasts} \\
C1 & 각 점수의 실행 시간과 그래프 크기 & 표~\ref{tab:benchmark-runtime} \\
C2 & 같은 기간 안에서 각 점수가 자기 화살표 종류의 점검과 맞는 정도예요. 이 일치는 어느 정도 처음부터 생길 수밖에 없고, EndorseRank에서는 동점(점수가 똑같은 것)이 그 정도를 정해요 & 표~\ref{tab:alignment-hybrid}와~\ref{tab:tie-sensitivity} \\
C3 & EndorseRank와 AWP는 코호트의 순서를 다르게 매겨요. 주로 어떤 지갑이 허락 그래프에 들어 있는지 때문이에요 & 표~\ref{tab:tie-sensitivity}; 그림~\ref{fig:rank-scatter} \\
C4 & 점수를 매긴 기간 밖에서는, 두 기간 어느 쪽에서도 어떤 페이지랭크도 같은 종류 화살표의 원시 차수만큼 라벨을 잘 줄 세우지 못해요 & 표~\ref{tab:holdout-spenders}, \ref{tab:holdout-diff}, \ref{tab:holdout-receiving}와~\ref{tab:fresh-posthoc} \\
C5 & 두 종류의 화살표를 결합해도, 그 두 층을 따로 걷는 것을 이기지 못해요. 또 등록 재현(계획을 먼저 공개해 두고 새 자료로 다시 해 보기)은 송금 걷기보다 낫다는 탐색적(미리 정하지 않고 살펴본) 결과를 확인해 주지 않아요. 새 허락에서 나아진 것은 다시 나타나지만, 새 송금자에서 비열등성(정한 폭보다 더 나쁘지 않음)은 통과하지 못하고, AWP와 견주면 분명히 실패해요 & 표~\ref{tab:alignment-hybrid}, \ref{tab:tau-diff-hybrid}, \ref{tab:holdout-diff}와~\ref{tab:fresh-contrasts}; \ref{sub:fresh-results}절 \\
\bottomrule
\end{tabular}
}
\end{table}

\dissertationSubheading[sec:research-objectives]{연구 목표}

이 연구에는 목표(O)가 다섯 개 있어요. RQ1--RQ5마다 하나씩이에요. 목표마다 그것을 실제로 잴 수 있게 만든 측정 방법과, \ref{ch:discussion}장에서 목표를 이루었는지 판단할 때 쓰는 기준을 함께 적었어요.

\begin{enumerate}
\item[O1.] 엔도스랭크와 AWP가 같은 코호트의 지갑들을 다른 순서로 세우는지 밝혀요(RQ1). 측정: $n=5{,}521$개 지갑에 대한 두 점수 벡터(점수를 차례로 늘어놓은 것) 사이의 켄달 $\tau$와 95\% 부트스트랩(다시 뽑기) 구간이에요. 기준: 구간에 $0$(서로 관계없는 순서)도 $1$(똑같은 순서)도 들어 있지 않아야 해요.
\item[O2.] 같은 기간 안에서 각 점수가 각 구성개념과 얼마나 잘 맞는지 따져요(RQ2). 측정: 모든 점수와 구성개념에 대해, 지표 가족(비슷한 지표 묶음) 단위의 켄달 $\tau$와 95\% 구간을 구해요. 기준: 어떤 지표 가족에서 점수의 구간이 $0$을 포함하지 않으면, 그 점수는 그 구성개념과 맞는다고 봐요. 점수가 자기 화살표의 구성개념과 맞는 것은 의도한 구성개념 점검으로 읽어요. 그 대리 지표들이 그래프와 같은 이벤트에서 나오기 때문이에요.
\item[O3.] 같은 지갑들에서 각 점수를 계산하는 데 얼마나 드는지 재요. 그리고 뽑은 지갑 수에 따라 실행 시간이 어떻게 바뀌는지 재요(RQ3). 측정: 푸는 시간(몸풀기 실행 한 번 뒤에 시간을 잰 실행 다섯 번의 평균과 표준편차), 메모리, 반복(한 바퀴) 횟수, $|V|$와 $|E|$를 모두 한 가지 절차로 재요. 기준: 시간 비율은 반드시 $|E|$의 비율과 함께 보고해요. 그래야 어떤 차이든 그래프 크기로 거슬러 올라가 설명할 수 있어요.
\item[O4.] $t_1$에 고정한 점수가 코호트의 스펜더들 사이에서 $t_2$의 새 허락과 새 송금자를 예측하는지 시험해요. 그리고 어떤 페이지랭크라도 자기가 다듬는 원시 차수보다 나은지 시험해요(RQ4). 측정: 2026년 2월 28일에 계산한 모든 점수(두 원시 차수도 포함)와, $n=1{,}335$개 스펜더에서 2026년 3월부터 5월까지 센 두 라벨 사이의 켄달 $\tau$와 95\% 구간이에요. 여기에 각 페이지랭크를 자기 차수와 맞대어 본 짝지은 대비도 더해요. 기준: 구간이 $0$을 포함하지 않으면 그 점수는 라벨을 예측한다고 봐요. 그리고 짝지은 차이의 구간이 $0$보다 위에 있을 때만 페이지랭크가 자기 차수보다 나아졌다고 봐요.
\item[O5.] 두 종류의 화살표를 가로지르는 걷기 하나가 송금만 걷는 걷기보다 나은지 시험해요. 그리고 어느 한 층만 걷는 걷기를 이기는지 시험해요(RQ5). 측정: $\lambda=0.5$의 C-PR과 S-PR을 써요. C-PR에는 민감도 분석(조건을 바꿔 결과가 얼마나 움직이는지 보기)을 위해 $0.25$와 $0.75$ 값도 써요. 이 점수들을 한 층만 걸은 두 걷기, 그리고 엔도스랭크, AWP와 견주어요. 견주기는 홀드아웃 라벨과 같은 기간 라벨에서 $\Delta\tau$에 대한 짝지은 구간으로 해요. 기준(각 기준은 그것으로 판단할 자료보다 먼저 정했어요): (a) 2026년 9월 14일의 규칙이에요. 이 규칙에 따르면 C-PR은 봄 홀드아웃에서 새 허락으로는 허락 걷기보다 나쁘지 않아야 하고, 새 송금자로는 송금 걷기보다 나아야 해요. 또 같은 기간 안에서는 송금 지표 가족과 허락 지표 가족에서 앞선 층의 구간 안에 들어야 하고, 시빌 안정성(가짜 주소 문제를 고려해 활동이 얼마나 꾸준한지 잰 것)에서는 두 층보다 모두 위에 있어야 해요. (b) 2026년 10월 1일에 등록한 규칙이에요. 이 규칙에 따르면 C-PR은 6월부터 8월까지의 기간에 새 허락 쌍에서 송금 걷기를 이겨야 해요. 그리고 새 송금자나 트레이더의 청산 없이 닫은 비율에서 0.02보다 더 크게 뒤지면 안 돼요. 미리 등록한 민감도 분석 하나는 송금 걷기 자리에 AWP를 넣고 (b)를 되풀이해요. 이 분석은 규칙을 바꾸지 않아요.
\end{enumerate}

\dissertationSubheading[sec:research-questions]{연구 질문}

처음 세 연구 질문에는 가설 세 개가 따라와요. H1: 같은 지갑 모음에서 엔도스랭크와 AWP는 서로 다른 지갑 순서를 만들어요. H2: 각 점수는 자기가 걷는 화살표로 만든 점검과 맞아요. H3: 같은 $n$에서, 점수를 계산하는 시간은 그 그래프의 화살표 수를 따라가요. 그래서 가장 듬성듬성한 그래프를 쓰는 엔도스랭크가 가장 빨리 계산돼요. RQ4와 RQ5는 관련 자료보다 먼저 정해 둔 대비와 규칙으로 답해요(\ref{sec:holdout-protocol}절). 표~\ref{tab:hrq-claim-map}에 질문들을 정리했어요.

\begin{enumerate}
\item[RQ1.] 같은 아비트럼 지갑 모음에서, 엔도스랭크와 AWP는 서로 다른 지갑 순서를 만들까요?
\item[RQ2.] 각 점수는 지금의 허락 점검, 그리고 지금의 송금 점검과 얼마나 잘 맞을까요(스피어먼의 $\rho$와 켄달의 $\tau$)?
\item[RQ3.] 같은 지갑 모음에서 각 점수를 계산하는 데, 시간, 메모리, 반복 횟수로 재면 얼마나 들까요?
\item[RQ4.] 시점 $t_1$에 계산한 점수가 있을 때, 그 점수는 $t_2$에 허락을 받는 스펜더에게 새로 생기는 주인--스펜더 허락과 새 송금자를 예측할까요? 그리고 어떤 페이지랭크라도 $t_1$에서 자기가 다듬는 원시 차수보다 그것들을 더 잘 예측할까요?
\item[RQ5.] 허락 층과 송금 층을 함께 걷는 결합 페이지랭크는 AWP처럼 송금만 걷는 걷기보다 나을까요? 그리고 어느 한 층만 걷는 걷기보다 지갑 순위를 더 잘 매길까요? 표본 안에서뿐 아니라 표본 밖에서도 그럴까요?
\end{enumerate}

RQ1은 한 층짜리 점수 두 개의 순서를 비교해요. RQ2에서는 자기 구성개념의 대리 지표가 모두 그래프와 같은 이벤트에서 계산돼요. 그래서 RQ2는 구성 타당도에 대한 관찰이지, 외부 검증(바깥 자료로도 맞는지 보기)이 아니에요. RQ3은 같은 지갑 모음에서 비용을 재고, 점의 개수에 따라 비용이 어떻게 바뀌는지 재요. RQ4는 보증을 구성개념으로 두고, 앞으로 생길 새 허락을 예측하는지 시험해요. 새 송금자는 두 번째 라벨로 써요. $t_1$의 원시 차수는 각 페이지랭크가 이겨야 할 기준선이에요. RQ4는 연구 계획서에 있던 거래 성공 검증을 대신해요. 그 검증에서 결과를 바탕으로 만든 비교 기준은 순환적(맞혀야 할 값으로 만든 기준)이었어요(\ref{sub:gf-pr}절). RQ5는 두 종류의 화살표를 비교하는 데서 나아가 둘을 합쳐 봐요. 두 층을 모두 이기는지에 대한 규칙은 결합 점수를 하나라도 계산하기 전에 정했어요. 송금 걷기보다 나은지에 대한 규칙은 봄 라벨을 본 뒤에 나왔어요. 그래서 그 규칙은 2026년 10월 1일에 등록한 6월부터 8월까지의 기간에서만 판단해요(\ref{sec:holdout-protocol}절과 \ref{sub:fresh-holdout}절).

\dissertationSubheading[sec:literature-summary]{문헌 검토 요약}

\ref{ch:literature}장은 네 분야를 살펴봐요. 그래프 순위 분야에서, 페이지랭크\cend{15}는 되풀이해 퍼뜨린 보증으로 점에 순위를 매겨요. 반복 한 번에 드는 비용은 화살표 수에 비례해요\cend{23}. EigenTrust\cend{16}와 TrustRank\cend{17}는 정규화(값을 일정한 기준에 맞추기)와 씨앗 집합을 써서 이 방법을 P2P 신뢰와 웹 스팸에 맞게 고쳤어요. 온체인 평판 분야에서, Do와 Do\cend{4}는 이더리움 거래에 페이지랭크와 HodgeRank를 썼어요. AWP\cend{3}는 거래 화살표에 금액과 지난 시간으로 무게를 달아요. 그리고 들어온 송금의 수와 금액으로 순위를 검증하는데, 이 통계들은 AWP의 그래프와 같은 자료에서 나온 것이에요. Lin 등\cend{12}은 거래 네트워크를 따라 계정의 위험을 퍼뜨려요. 블록체인 평판 시스템을 정리한 조사 연구들\cend{25,26}은 분명한 평가와 송금 양이 흔히 쓰는 신뢰 신호라는 것을 알아냈어요. 디파이에서 실제로 쓰는 도구는 또 다르게 작동해요. 휴먼 패스포트(Human Passport)는 주소 주인에 대한 온체인과 오프체인의 자격 증명(그 사람에 대해 확인해 준 사실)을 모아, 그 주소가 다른 사람과 구별되는 한 사람의 것인지 판단해요\cend{64}. 이 연구들 가운데 어느 것도 ERC-20 허락\cend{18}을 그래프의 화살표로 쓰지 않아요. 디파이 위험과 측정 분야에서, 조사 연구들\cend{8,27}은 거래 상대의 행동 위험을 계약, 오라클(블록체인 밖의 값을 들여오는 장치), 시장, 거버넌스(운영 규칙을 정하는 일) 위험과 나란히 놓아요. 그리고 자전 거래(자기끼리 사고팔아 거래가 많은 척하기)\cend{28,29}, 추출 가능한 가치(거래 순서를 이용해 빼내는 이익)\cend{30}, 사기 토큰\cend{31}, 에어드롭 파밍(공짜 토큰을 노린 가짜 활동)\cend{32}은 송금 활동의 많은 부분이 꾸며 낸 것임을 보여 줘요. 측정 이론 분야에서, 구성 타당도\cend{22,33}는 점수가 무엇을 재는지를 바깥 기준과 구별해요. 그리고 부트스트랩\cend{34,35}은 분포를 미리 가정하지 않고 순위 통계의 구간을 구해 줘요. 빈틈(\ref{sec:research-gap}절)은 바로 이것이에요. 허락 화살표로 매긴 순위를 하나의 코호트에서 송금 방법, 원시 차수와 나란히 평가하는 것이에요. 이때 신뢰구간, 짝지은 차이, 정해 둔 시간 재기 절차, 시간 홀드아웃을 함께 써요.

\dissertationSubheading[sec:methodology-summary]{연구 방법 요약}

\ref{ch:methodology}장은 공개 장부 자료에 수를 세어 비교하는 설계(양적 비교 설계)를 써요. 2025년 12월 1일부터 2026년 5월 31일까지 아비트럼 원의 ERC-20 \texttt{Approval} 로그와 \texttt{Transfer} 로그를 구글 빅쿼리(Google BigQuery)에서 꺼내요. 대상은 GMX~V2 지갑 $n=5{,}521$개로 된 매칭 코호트예요. 이 코호트는 GMX 자료만으로 정했어요. 거듭제곱 반복(같은 계산을 되풀이하기)으로 페이지랭크를 계산하는 프로그램 하나가 마지막 허락 그래프와 송금 그래프에 점수를 매겨요. 결합 연산자도 같은 풀이 프로그램으로 돌려요. 각 점수는 스피어먼의 $\rho$와 켄달의 $\tau$로 비교해요. 같은 기간 안에서는 송금, 허락, 시빌 안정성 대리 지표와 견주어요. 기간 밖에서는 스펜더 코호트의 새 허락, 새 송금자와 견주고, 동결일(점수를 고정하는 날)의 원시 차수를 기준선으로 써요. 미리 등록한 두 번째 기간은 2026년 6월부터 8월까지 홀드아웃을 되풀이하고, 트레이더 라벨을 하나 더해요. 모든 계수(일치 정도를 나타내는 수)에는 짝지은 지갑 다시 뽑기 400번(시드 42)으로 구한 95\% 구간이 붙어요. 방법끼리의 비교도 짝지어서 해요. 실행 시간, 메모리, 반복 횟수, 그래프 크기는 한 가지 시간 재기 절차로 기록해요. 모든 표 뒤에 있는 코드, 설정, 기계가 읽을 수 있는 요약은 공개 저장소에 있어요(\ref{sec:reproducibility}절). 그리고 이 연구는 공개 자료만 써요(\ref{sec:ethics}절).

\dissertationSubheading[sec:contributions]{기여}

이 논문은 온체인 평판과 디파이 위험 연구에 네 가지 기여(이 연구가 새로 보탠 것)를 해요. 엔도스랭크와 AWP는 둘 다 무게를 단 페이지랭크\cend{15}를 한 종류의 화살표에 써요. 둘은 저마다 자기 재시작 분포(어디서 다시 출발하는지의 비율)를 가져요. 이동 규칙을 바꾸는 것은 결합 연산자뿐이에요. 그러므로 새로운 것은 화살표 종류 하나, 평가 설계 하나, 그리고 두 종류의 화살표가 합쳐지는지에 대한 등록된 시험 하나예요.
\begin{itemize}
\item 마지막 허락으로 화살표 만들기(엔도스랭크): 보증 그래프를 만들어요. 이 그래프에서 화살표 $u \to v$는 스펜더 $v$가 주인 $u$에게서 받아 둔, 기간 안의 마지막 ERC-20 허락을 나타내요. 이 허락은 \texttt{approve}나 \texttt{permit}으로 정해져요. 화살표 무게는 AWP가 송금에 무게를 다는 방식과 똑같이 달아요. 이 화살표는 돈을 쓸 권리를 맡긴 것(위임)을 기록하지, 돈을 낸 기록이 아니에요. 그리고 살펴본 평판 방법 가운데 이 화살표를 쓰는 것은 하나도 없어요(\ref{sec:research-gap}절).
\item 통계적 추론(표본으로 전체를 짐작하기)을 갖춘 같은 코호트 평가: 이 논문의 모든 점수는 한 지갑 모음에서 계산해요. 그리고 AWP 검증 절차\cend{3}에 따라 똑같은 송금, 허락, 시빌 안정성 대리 지표로 점검해요. 모든 계수에는 95\% 부트스트랩 구간을, 대비에는 짝지은 구간을 붙여요\cend{34}. 각 점수가 자기 대리 지표와 맞는 것의 일부는 처음부터 생길 수밖에 없어요. 그리고 실행 시간은 화살표 수와 나란히 보고해요. 주장~C1은 희소성(화살표가 적음)에서 나오기 때문이에요.
\item 차수 기준선을 갖춘 시간 홀드아웃: $t_1$에 고정한 점수를 $t_1$ 뒤에 생긴 새 허락과 새 송금자로 시험해요. 이때 페이지랭크가 다듬는 원시 차수도 나란히 시험해요. 이 설계는 표본 안의 일치와 표본 밖의 예측을 떼어 놓고, 퍼뜨린 점수와 그 자리에서 센 수도 떼어 놓아요. 공개된 처리 과정(파이프라인, 부록~\ref{ch:appendix-eval})과 함께, 이 설계는 주장~C4와~C5를 뒷받침하는 기간 밖 증거를 줘요.
\item 결합 연산자와 그 등록된 평가: C-PR은 금액을 그대로 쓴 허락 층과 금액을 그대로 쓴 송금 층을 점마다 정해진 비율 $\lambda$로 섞어요. S-PR은 씨앗을 쓰는 빼 보기 실험이에요. 첫 번째 규칙은 결합 점수를 하나라도 계산하기 전에 $\lambda$와 함께 정했어요. 이 규칙은 C-PR이 한 층만 걷는 두 걷기를 모두 이기기를 요구했어요. 두 번째 규칙은 그 기간의 로그를 꺼내기 전에 등록했어요. 이 규칙은 C-PR이 송금 걷기보다 나아지기를 요구했어요. 두 규칙 모두 통과하지 못했고, 이 논문은 실패했다고 그대로 보고해요(주장~C5). 빼 보기 실험은 결과를 하나 더 보태요. 허락 화살표가 걷기를 실어 나를 때는 허락 신호가 살아남아요. 하지만 허락 화살표가 걷기가 다시 출발할 곳만 정할 때는 그 신호가 사라져요.
\end{itemize}


\dissertationSubheading[sec:scope-limitations]{범위와 한계}

이 연구는 2025-12-01부터 2026-05-31까지의 아비트럼 원을 다뤄요. 주 표본은 그 기간에 포지션을 세 번 이상 닫은 GMX~V2 지갑 $n=5{,}521$개예요. 홀드아웃 표본은 2026년 2월 28일에 0보다 큰 허락을 가지고 있던 스펜더 $n=1{,}335$개예요. 엔도스랭크는 지갑이 받은 허락으로 점수를 매겨요. 그래서 스펜더들은 가려내지만, 허락을 해 주는 주인들은 가려내지 못해요. 매칭된 트레이더 대부분은 주인이에요. 그래서 홀드아웃은 스펜더를 대상으로 하고, 미리 등록한 6월부터 8월까지의 기간은 트레이더 라벨을 하나 더해요(\ref{sub:fresh-holdout}절). 실행 시간은 코호트 그래프의 부분 그래프에서도 재요. 이 부분 그래프는 더 큰 주소 풀($N=99{,}080$)을 거쳐 뽑지만, 그 풀은 화살표를 하나도 더하지 않아요. 담보 없는 대출의 채무 불이행(빌린 돈을 갚지 못함) 라벨은 이 연구의 범위 밖이에요. 더 넓힌 기준선들은 따로 보관해 두었어요(부록~\ref{ch:appendix-archive}). 그리고 결과가 아비트럼 원이나 살펴본 달들 밖에서도 들어맞는다고 주장하지 않아요. 금액은 토큰의 원래 단위(소수점 조정을 하지 않은 가장 작은 단위) 그대로예요. 엔도스랭크와 AWP는 금액을 상한이 있는 함수에 넣어요. 그래서 원래 단위로 10개 이상인 허락이나 송금은 모두 거의 똑같이 쳐요. 반면 C-PR의 금액을 그대로 쓴 층에서는 무제한 허락(쓸 수 있는 양에 끝이 없는 허락)이 많은 주인의 줄을 좌우해요(\ref{ch:methodology}장). 허락이나 송금을 나쁜 뜻으로 조작하는 일은 실제 위협이에요. \ref{ch:discussion}장은 시빌 주소(가짜로 만든 주소)로 이루어진 닫힌 농장(공격자가 만든 가짜 주소 무리)이 엔도스랭크를 싸게 부풀릴 수 있음을 보여 줘요. 그리고 시빌 보정(가짜 주소 문제를 고려해 고친) 대리 지표들은 몇 가지 조작 방식을 설명할 뿐, 점수를 그 조작에서 지켜 주지는 못해요\cend{11}.

\dissertationSubheading[sec:roadmap]{연구 길잡이}

\ref{ch:literature}장은 그래프를 바탕으로 한 온체인 평판, ERC-20 허락의 뜻, 디파이에서 실제로 쓰는 평판 도구를 살펴보고, 이론 틀을 세워요. \ref{ch:methodology}장은 코호트, 그래프와 점수, 구성개념 점검, 시간 홀드아웃, 시간 재기 절차, 그리고 재현성(다른 사람이 똑같이 다시 해 볼 수 있음)과 윤리의 기준을 정해요. \ref{ch:results}장은 속도 재기인 벤치마크(주장~C1), 모든 점수의 같은 기간 일치도(C2), 순위 갈라짐(C3), 강건성 점검(조건을 바꿔도 결과가 버티는지 보기), 그리고 등록 재현을 포함한 시간 홀드아웃(C4와 C5)을 보고해요. 부록~\ref{ch:appendix-er-awp}에는 엔도스랭크와 AWP를 바로 비교한 결과를 모았어요. \ref{ch:discussion}장은 연구 질문과 목표에 따라 결과를 풀이하고, 그 결과가 무엇을 뜻하는지 따져요. 또 허락 그래프에 대한 시빌 공격(가짜 주소를 많이 만들어 속이는 공격)의 비용을 모형으로 만들고, 타당성을 위협하는 것과 한계를 따져 봐요. \ref{ch:conclusion}장은 RQ1--RQ5에 답하고, 지식에 대한 기여를 밝히고, 앞으로 할 일을 내놓아요.

